\subsection{利普希茨方程与局部利普希茨方程的解的存在性与唯一性}

定理 1（柯西）. 设 \(\mathbf{f}\) 是 \(\mathrm{I} \times \mathrm{H}\) 上的利普希茨函数，设 \(\mathrm{J}\) 是 \(\mathrm{I}\) 的一个不退化为一点的紧子区间，\(t_0\) 是 \(\mathrm{J}\) 的一点，\(\mathrm{S}\) 是以 \(\mathbf{x}_0\) 为中心、以 \(r\) 为半径、包含于 \(\mathrm{H}\) 中的开球，而 \(\mathrm{M}\) 是 \(\| \mathbf{f}(t, \mathbf{x})\|\) 在 \(\mathrm{J} \times \mathrm{S}\) 上的上确界。在这些条件下，对每个紧区间 \(\mathrm{K}\)，若它不退化为一点、包含于 \(\mathrm{J}\) 与 \(]t_0 - r / M\)、\(t_0 + r / M[\) 的交中且包含 \(t_0\)，则存在微分方程 \(\mathbf{x}' = \mathbf{f}(t, \mathbf{x})\) 的唯一一个解，它定义在 \(\mathrm{K}\) 上、取值于 \(\mathrm{S}\) 中、且在点 \(\mathbf{x}_0\) 处等于 \(t_0\) 。

的确，当 \(\varepsilon > 0\) 足够小时，\(F_{\varepsilon}\)（即 (1) 的精确到 \(\varepsilon\) 、定义在 K 上、取值于 S 中且等于 \(x_{0}\) 于点 \(t_{0}\) 的近似解）非空（IV，第 166 页，命题 3）；此外，若 u 与 v 都属于 \(F_{\varepsilon}\) ，则由 (15)（IV，第 170 页）

\[
\| \mathbf {u} (t) - \mathbf {v} (t) \| \leqslant 2 \varepsilon \frac {e ^ {k | t - t _ {0} |} - 1}{k}
\]

对一切 \(t \in K\) ，于是诸集 \(F_{\varepsilon}\) 构成一个滤基 G，它在 K 上一致收敛于一个连续函数 w，后者等于 \(x_{0}\) 于 \(t_{0}\) ；而且 w 取值于 S 中，因为当 \(\varepsilon\) 足够小时，诸函数 \(u \in F_{\varepsilon}\) 的值都属于一个包含于 S 中的闭球。由于 \(\mathbf{f}(t, \mathbf{u}(t))\) 沿 G 在 K 上一致趋于 \(\mathbf{f}(t, \mathbf{w}(t))\) ，w 满足 IV 第 165 页的方程 (6)，因而是 (1) 的解。解的唯一性由 IV 第 170 页的不等式 (15) 立即得出，其中取 \(\varepsilon_{1} = \varepsilon_{2} = 0\) 与 \(\mathbf{u}(t_{0}) = \mathbf{v}(t_{0})\) 。

我们说函数 \(\mathbf{f}\)（定义在 \(\mathrm{I} \times \mathrm{H}\) 上）是局部利普希茨的，如果对 \((t, \mathbf{x})\) 的每个点 \(\mathrm{I} \times \mathrm{H}\) ，存在 \(\mathrm{V}\) 的一个邻域 \(t\)（关于 \(\mathrm{I}\) ）与 \(\mathrm{S}\) 的一个邻域 \(\mathbf{x}\)，使得 \(\mathbf{f}\) 在 \(\mathrm{V} \times \mathrm{S}\) 上是利普希茨的（常数 \(k\) 依赖于 \(\mathrm{V}\) 与 \(\mathrm{S}\) ）。由博雷尔–勒贝格定理，对每个紧区间 \(\mathrm{J} \subset \mathrm{I}\) 与每个点 \(\mathbf{x}_0 \in \mathrm{H}\)，存在以 \(\mathrm{S}\) 为中心、包含于 \(\mathbf{x}_0\) 中的开球 \(\mathrm{H}\)，使得 \(\mathbf{f}\) 在 \(\mathrm{J} \times \mathrm{S}\) 上是利普希茨的；于是 \(\mathbf{f}\) 满足 IV 第 3 页引理 1 的假设。当 \(\mathbf{f}\) 在 \(\mathrm{I} \times \mathrm{H}\) 上是局部利普希茨的时候，我们称方程 \(\mathbf{x}' = \mathbf{f}(t, \mathbf{x})\) 在 \(\mathrm{I} \times \mathrm{H}\) 上是局部利普希茨的。

我们要对局部利普希茨方程推广并阐明 IV 第 10 页的定理 1；我们限于 \(t_0\) 是区间 I 的左端点的情形；对 \(t_0\) 是 I 的任意点的情形，容易化归为上述情形（参看 IV，§ IV ??，第 9 页，推论）。

定理 2. 设 \(\mathbf{I} \subset \mathbf{R}\) 是一个（不退化为一点的）区间，其左端点为 \(t_0 \in \mathbf{I}\) ；设 \(\mathbf{H}\) 是 \(\mathbf{E}\) 中的一个非空开集，并设 \(\mathbf{f}\) 是 \(\mathbf{I} \times \mathbf{H}\) 上的局部利普希茨函数。

\(1^{\circ}\) 对 \(\mathbf{x}_0\in \mathrm{H}\) ，存在最大的区间 \(\mathbf{J}\subset \mathbf{I}\)（其左端点为 \(t_0\in \mathbf{J}\) ），在该区间上存在一个积分 \(\mathbf{u}\) ，即方程 \(\mathbf{x}' = \mathbf{f}(t,\mathbf{x})\) 的、取值于 H 中且等于 \(\mathbf{x}_0\) 于点 \(t_0\) 的积分；这个积分是唯一的。

\(2^{\circ}\) 若 \(\mathrm{J} \neq \mathrm{I}\) ，则 \(\mathrm{J}\) 是长度有限的半开区间 \([t_0, \beta[\) ；此外，对每个紧子集 \(\mathbf{K} \subset \mathbf{H}\) ，集合 \(\overline{\mathbf{u}}^{-1}(\mathbf{K})\) 是 \(\mathbf{R}\) 的紧子集。

\(3^{\circ}\) 若 \(\mathbf{J}\) 有界，且若 \(\mathbf{f}(t,\mathbf{u}(t))\) 在 \(\mathbf{J}\) 上有界，则 \(\mathbf{u}(t)\) 在 \(\mathbf{c}\) 的右端点处有左极限 \(\mathbf{J}\) ；此外，若 \(\mathbf{J} \neq \mathbf{I}\) ，则 \(\mathbf{c}\) 是 \(\mathbf{H}\) 在 \(\mathbf{E}\) 中的边界点。

\(1^{\circ}\) 设 \(\mathfrak{M}\) 是由这样的区间 L（不退化为一点、左端点为 \(t_0 \in L\) 、包含于 I）组成的集合：在 L 上存在 (1)（IV，第 163 页）的、取值于 H 中且等于 \(\mathbf{x}_0\) 于 \(t_0\) 的解；由定理 1（IV，第 171 页），集合 \(\mathfrak{M}\) 非空。设 L 与 L' 是属于 \(\mathfrak{M}\) 的两个区间，例如设 L \(\subset\) L'；若 u 与 v 是 (1) 的两个积分，分别定义在 L 与 L' 上、取值于 H 中且等于 \(\mathbf{x}_0\) 于 \(t_0\) ，我们将看到 v 是 u 的扩张。的确，设 \(t_1\) 是使 \(t \in L\) 满足 \(\mathbf{u}(s) = \mathbf{v}(s)\) （对 \(t_0 \leqslant s \leqslant t\) ）的 t 组成的集合的上确界；我们将证明 \(t_1\) 是 L 的右端点。若非如此，则由连续性有 \(\mathbf{u}(t_1) = \mathbf{v}(t_1)\) ，且 \(\mathbf{x}_1 = \mathbf{u}(t_1)\) 属于 H；由于 f 是局部利普希茨的，定理 1 表明只能存在 (1) 的一个在 \(t_1\) 的邻域上定义、取值于 H 中且等于 \(\mathbf{x}_1\) 于 \(t_1\) 的积分；因此假定 \(t_1\) 不是 L 的右端点便导致矛盾。现在我们看到，若 J 是诸区间 L \(\in\) M 的并，则存在 (1) 的唯一一个积分 u，它定义在 J 上、取值于 H 中且等于 \(\mathbf{x}_0\) 于 \(t_0\) 。

\(2^{\circ}\) 设 \(\mathrm{J} \neq \mathrm{I}\) ，并设 \(\beta\) 是 \(\mathrm{J}\) 的右端点；若 \(\beta\) 是 \(\mathrm{I}\) 的右端点，则 \(\beta \in \mathrm{I}\) （于是 \(\beta\) 是有限的），且由假设 \(\mathrm{J} = [t_0, \beta[\) 。今设 \(\beta\) 不是 \(\mathrm{I}\) 的右端点；若 \(\beta \in \mathrm{J}\) ，则 \(\mathbf{u}(\beta) = \mathbf{c}\) 属于 \(\mathrm{H}\) ；由定理 1，存在 (1) 的一个取值于 \(\mathrm{H}\) 中、定义在一个区间上的积分

\[
[ \beta , \beta_ {1} [ \subset I
\]

且等于 \(\mathbf{c}\) 于 \(\beta\) ；这样 \(\mathrm{J}\) 就不是 \(\mathfrak{M}\) 中最大的区间，这是荒谬的；故 \(\mathrm{J} = [t_0, \beta[\) 。

若 K 是 H 的紧子集，则 \(\overline{\mathbf{u}}^{-1}(\mathbf{K})\) 在 J 中是闭的；我们将看到存在 \(\gamma\in J\) 使得 \(\overline{\mathbf{u}}^{-1}(\mathbf{K})\) 包含于 \([t_{0},\gamma]\) 中，这就证明 \(\overline{\mathbf{u}}^{-1}(\mathbf{K})\) 是紧的。若不然，就会存在一点 \(c\in K\) 使得 \((\beta,\mathbf{c})\) 是点 \((t,\mathbf{u}(t))\)（其中 \(t<\beta\) 与 \(\mathbf{u}(t)\in\mathbf{K}\) ）所成之集的聚点。由于 \(\beta\in I\) 与 \(c\in H\) ，则存在 \(\beta\) 在 I 中的邻域 V，以及以 c 为中心、半径 r 且包含于 H 中的开球 S，使得 f 在 V×S 上是利普希茨的且有界；令 M 为 \(\|\mathbf{f}(t,\mathbf{x})\|\) 在该集合上的上确界。由假设存在 \(t_{1}\in J\) 使得 \(\beta-t_{1}<r/2M\) 、\(t_{1}\in V\) 与 \(\|\mathbf{u}(t_{1})-\mathbf{c}\|\leqslant r/2\) ；定理 1 表明存在 (1) 的唯一一个积分，它取值于 H 中、定义在包含 \([t_{1},t_{2}]\) 的区间 \(\beta\) 上、且等于 \(\mathbf{u}(t_{1})\) 于 \(t_{1}\) ；由于这个区间在区间 \([t_{1},\beta[\) 上与 u 一致，便推出 \(J=[t_{0},\beta[\) 不是 M 中最大的区间，这是荒谬的。

\(3^{\circ}\) 设 \(\mathbf{J}\) 有界且 \(\| \mathbf{f}(t,\mathbf{u}(t))\| \leqslant M\) 在 \(\mathbf{J}\) 上成立；则在 \(\| \mathbf{u}'(t)\| \leqslant \mathbf{M}\) 于 \(\mathbf{J}\) 的一个可数子集的补集上成立；于是由中值定理，对 \(\| \mathbf{u}(s) - \mathbf{u}(t)\| \leqslant \mathbf{M} |s - t|\) 与 \(s\)（都属于 \(t\) ）有 \(\mathbf{J}\) ；由柯西准则，\(\mathbf{u}\) 有左极限 \(\mathbf{c}\) 于 \(\beta\) 处（它是 \(\mathbf{J}\) 的右端点）；若 \(\mathbf{J} \neq \mathbf{I}\) ，则 \(\mathbf{c}\) 不可能属于 \(\mathbf{H}\) ，因为把 \(\mathbf{u}\) 在 \(\beta\) 处连续延拓后，\(\mathbf{u}\) 就成为 (1) 的一个取值于 \(\mathbf{H}\) 中、定义在区间 \([t_0, \beta]\) 上且等于 \(\mathbf{x}_0\) 于 \(t_0\) 的积分；于是就有 \(\mathbf{J} = [t_0, \beta]\) ，这与我们在 \(2^{\circ}\) 中已看到的相矛盾。

推论 1. 若 H = E 且 J ≠ I，则 \(\mathbf{f}(t, \mathbf{u}(t))\) 在 J 上不是有界的；进而，若 E 是有限维的，则 \(\| \mathbf{u}(t) \|\) 在 J 的右端点处有左极限 \(+\infty\) 。

第一部分是定理 2 第三部分的直接推论。若 E 是有限维的，则每个闭球 \(S \subset E\) 是紧的，于是定理 2 的第二部分表明存在 \(\gamma \in J\) 使得 \(\mathbf{u}(t) \notin \mathbf{S}\) 对 \(t > \gamma\) 成立。

若 E 是有限维的，可能发生这样的情况：\(J \neq I\) 而 \(\|\mathbf{u}(t)\|\) 当 t 趋于 J 的右端点时保持有界（IV，第 200 页，习题 5）。

推论 2. 若在 I × H 上函数 f 关于一个规则函数 k(t) 是利普希茨的，且若 J 的右端点 β 属于 I，则 u 在 β 处有左极限；若 H = E 且 f 在 I × E 上关于一个规则函数 k(t) 是利普希茨的，则 J = I。

的确，若 \(\beta \in \mathbf{I}\) ，则存在 \(\mathrm{V}\) 在 \(\beta\) 中的一个紧邻域 \(\mathbf{I}\) ，使得 \(\mathbf{f}(t, \mathbf{x}_0)\) 与 \(k(t)\) 在 \(\mathrm{V}\) 上有界；于是 \(\| \mathbf{f}(t, \mathbf{x}) \| \leqslant m \| x \| + h\) （\(m\) 与 \(h\) 为常数）

在 \(V \times H\) 上成立，由此 \(\left\|\mathbf{u}^{\prime}(t)\right\| \leqslant m \left\|\mathbf{u}(t)\right\| + h\) 在 \(V \cap J\) 的一个可数子集的补集上成立，从而 \(\left\|\mathbf{u}(t)\right\| \leqslant m \int_{t_{0}}^{t} \left\|\mathbf{u}(s)\right\| ds + q\) （q 为常数）在 \(V \cap J\) 上成立；引理 2（IV，第 170 页）表明 \(\left\|\mathbf{u}(t)\right\| \leqslant c e^{mt} + d\) （c 与 d 为常数）在 \(V \cap J\) 上成立，于是 \(\mathbf{f}(t, \mathbf{u}(t))\) 在 J 上保持有界，从而由 IV 第 172 页的定理 2 即得本推论。

例. 1) 对形如 \(\mathbf{x}' = \mathbf{g}(t)\) 的微分方程（其中 \(\mathbf{g}\) 在 I 上是规则的），每个积分 \(\mathbf{u}\) 显然都定义在整个 I 上。应当注意，\(\mathbf{u}\) 可以在 I 上有界而 \(\mathbf{g}(t)\) 并不如此。

\begin{enumerate}
\def\labelenumi{\arabic{enumi})}
\setcounter{enumi}{1}
\item
  对数量方程 \(x' = \sqrt{1 - x^2}\) 有 \(I = R\) ，\(H = ] - 1, +1[\) 。若取 \(t_0 = x_0 = 0\) ，则相应的积分是 \(\sin t\) ，其定义区间是包含 0 且使 \(\sin t\) 的导数为正的最大区间，也就是说，\(] - \pi / 2, +\pi / 2[\) ；在这个区间的端点处，积分趋于 \(H\) 的一个端点。
\item
  对数量方程 \(x' = 1 + x^{2}\) 有 I = H = R；在 t = 0 处为零的积分是 \(\tan t\) ，而包含 0 且使该函数连续的最大区间是 \(J = ] - \pi/2, +\pi/2[\) ；且 \(|\tan t|\) 在 J 的端点处趋于 \(+\infty\) （参看 IV，第 173 页，推论 1）。
\item
  对数量方程 \(x' = \sin tx\) 有 I = H = R，且右端在 I × H 上有界，故（IV，第 173 页，推论 1）每个积分都定义在整个 R 上。
\end{enumerate}

